% Sustentacion del laboratorio de modelado supervisado
% Compilar desde la raiz:
% latexmk -pdf -interaction=nonstopmode -halt-on-error -outdir=build/slides slides/presentacion.tex

\documentclass[11pt]{beamer}

\usetheme{Madrid}
\usecolortheme{beaver}
\setbeamertemplate{navigation symbols}{}
\setbeamertemplate{caption}[numbered]

\usepackage[utf8]{inputenc}
\usepackage[T1]{fontenc}
\usepackage[spanish]{babel}
\usepackage{lmodern}
\usepackage{graphicx}
\usepackage{booktabs}
\usepackage{amsmath}

\graphicspath{{slides/}{reports/}{figures/modelado_proyecto/}{results/modelado_didactico/}{./}}

\makeatletter
\providecommand{\input@path}{}
\edef\input@path{{reports/tablas/}{./}}
\makeatother
\input{lab005_cifras}

\title[Modelado supervisado]{Modelado supervisado con regresion multiple}
\subtitle{Estimacion de la fraccion pulmonar de referencia a partir de la segmentacion automatica}
\author[Trujillo \and Martinez]{Juanita Trujillo Narvaez \\ Karol Isabella Martinez Villarreal}
\institute[ECI]{Ingenieria Biomedica \\ PSIM - 2026II - 80 \\ Docente: Ing. Pablo Eduardo Caicedo-Rodriguez, Ph.D.}
\date{\today}

\logo{\includegraphics[height=0.7cm]{Escuela_Rosario_logo.png}}

\begin{document}

% ---------------------------------------------------------------- 1
\begin{frame}[plain]
    \titlepage
\end{frame}

% ---------------------------------------------------------------- 2
\begin{frame}{Problema, salida y protocolo}
\footnotesize
\begin{columns}[T]
\begin{column}{0.50\textwidth}
\textbf{Pregunta}

?`Los descriptores de la mascara \emph{automatica} permiten estimar la
fraccion pulmonar de la mascara \emph{manual} en pacientes no usados para
ajustar el modelo?

\vspace{0.2cm}
\begin{block}{Contrato}
\scriptsize
\begin{tabular}{@{}ll@{}}
Unidad & paciente (\texttt{grupo\_id}) \\
Salida & \texttt{ref\_fraccion\_area} (adim.) \\
Referencia & mascara trazada por personas \\
Estrategia & \texttt{base\_adap}, elegibles \\
Metrica & RMSE, fijada antes \\
Ruta binaria & no aplica: TB es control \\
\end{tabular}
\end{block}
\end{column}

\begin{column}{0.48\textwidth}
\textbf{Del EDA al protocolo} (solo train)
\scriptsize
\begin{itemize}
    \item Predictores por asociacion con el objetivo y signo consistente
    entre origenes: \texttt{fraccion\_area} ($\rho=\ModRhoFA$),
    \texttt{energia\_wavelet\_d1} ($\rho=\ModRhoEW$),
    \texttt{simetria\_izq\_der} ($\rho=\ModRhoSI$)
    \item Excluidos: \texttt{ref\_}, \texttt{err\_}, Dice, Jaccard, IDs y
    origen
    \item Efecto de centro: fraccion media \ModRefMediaMC{} (MC) frente a
    \ModRefMediaSH{} (SH) $\Rightarrow$ errores por origen
\end{itemize}

\vspace{0.1cm}
\begin{tabular}{@{}lrrr@{}}
\toprule
 & Mascara & Elegibles & Pacientes \\
\midrule
Train & \ModMascTrain & \ModElegTrain & \ModNTrain \\
Valid. & \ModMascVal & \ModElegVal & \ModNVal \\
Prueba & \ModMascTest & \ModElegTest & \ModNTest \\
\bottomrule
\end{tabular}

\vspace{0.1cm}
La prueba se extrajo \textbf{despues} de congelar el protocolo.
\end{column}
\end{columns}
\end{frame}

% ---------------------------------------------------------------- 3
\begin{frame}{Modelos y seleccion}
\footnotesize
\begin{columns}[T]
\begin{column}{0.50\textwidth}
\textbf{Candidatos} ($z_j$ estandarizado con train)
\scriptsize
\begin{align*}
\text{G1:}\ \hat{y} &= \beta_0 + \textstyle\sum_j \beta_j z_j \\
\text{G2:}\ \hat{y} &= \text{G1} + \textstyle\sum_{j\le k} \beta_{jk} z_j z_k
\end{align*}
Imputacion, escalamiento y expansion dentro del \emph{pipeline}, ajustados
solo en train. Referencia: media de train.

\vspace{0.2cm}
\textbf{Practica comun (sinteticos)}
\begin{itemize}
    \item Continua: \DidModeloReg{} reconoce $x_0^2$; RMSE prueba
    \DidRMSETestReg{} vs \DidRMSETestBase{}
    \item Binaria: AUC \DidAUCTest{}, AP \DidAPTest{} (prevalencia
    \DidPrevTest{}); sensibilidad \DidSensTest{} con umbral 0.5
\end{itemize}
\end{column}

\begin{column}{0.48\textwidth}
\textbf{Proyecto: validacion} ($n=\ModNVal$)

\vspace{0.1cm}
\begin{tabular}{@{}lr@{}}
\toprule
Modelo & RMSE \\
\midrule
Media de train & \ModRMSEValBase \\
Grado 1 & \ModRMSEValUno \\
\textbf{Grado 2} & \textbf{\ModRMSEValDos} \\
\bottomrule
\end{tabular}

\vspace{0.2cm}
\scriptsize
Se elige \textbf{\ModElegidoNombre{}} por la regla predeclarada. Su ventaja
sobre grado 1 es modesta: casi todo el poder esta en el termino lineal.

\vspace{0.1cm}
Rango completo (\ModRangoDiseno), condicion \ModCondicion: coeficientes
identificables. \texttt{fraccion\_area} domina; las interacciones son
pequenas. Asociacion, no causalidad.
\end{column}
\end{columns}
\end{frame}

% ---------------------------------------------------------------- 4
\begin{frame}{Resultados en prueba y diagnostico}
\footnotesize
\begin{figure}
\centering
\includegraphics[width=0.86\textwidth]{prueba_prediccion_residuos.png}
\end{figure}

\vspace{-0.3cm}
\begin{columns}[T]
\begin{column}{0.48\textwidth}
\textbf{Evaluacion unica, $n=\ModNTest$ pacientes}
\scriptsize

RMSE \textbf{\ModRMSETest} vs \ModRMSETestBase{} de la media
(\ModMejoraTest\,\% menos), $R^2=\ModRTest$, igual a validacion.

SH: \ModRMSETestSH{} ($n=\ModNTestSH$) \quad MC: \ModRMSETestMC{}
($n=\ModNTestMC$). Residuo medio $\ModResTestSH$ y $\ModResTestMC$.
\end{column}

\begin{column}{0.48\textwidth}
\textbf{Limites}
\scriptsize
\begin{itemize}
    \item Regresion a la media en fracciones extremas
    \item Montgomery con pocos casos; sin tercer centro
    \item No aplica si la segmentacion falla (1 componente)
    \item Sin intervalos de incertidumbre
\end{itemize}
\end{column}
\end{columns}
\end{frame}

% ---------------------------------------------------------------- 5
\begin{frame}[fragile]{Reproduccion y siguiente decision}
\footnotesize
\begin{columns}[T]
\begin{column}{0.52\textwidth}
\textbf{Un comando reconstruye todo}
\scriptsize
\begin{verbatim}
uv sync --locked
uv run ruff check .
uv run pytest
uv run python scripts/reproduce.py
\end{verbatim}

\textbf{Comprobaciones automaticas}
\begin{itemize}
    \item Paciente que cruza particiones $\Rightarrow$ error
    \item Objetivo, IDs o \texttt{ref\_} como predictor $\Rightarrow$ error
    \item Imputacion aprendida solo en train
    \item Hashes de train y validacion verificados
    \item Prueba reconstruida: coeficientes y metricas identicos
\end{itemize}
\end{column}

\begin{column}{0.46\textwidth}
\textbf{Conclusion}
\scriptsize

Los descriptores automaticos estiman la fraccion pulmonar de referencia
con error de \ModRMSETest{} en pacientes nuevos, del mismo orden en ambos
centros.

\vspace{0.3cm}
\textbf{Siguiente decision}

Calcular intervalos por bootstrap de pacientes y comparar el error con la
variabilidad entre observadores, para decidir si la correccion es util en
la practica.
\end{column}
\end{columns}
\end{frame}

\end{document}